% LTeX: language=en-US
\documentclass[handout,usepdftitle=false,xcolor=dvipsnames,aspectratio=169]{beamer}
%\documentclass[beamer,usepdftitle=false,xcolor=dvipsnames,aspectratio=169]{beamer}
%\documentclass[trans,notes=show,usepdftitle=false,aspectratio=169]{beamer}
 

\usepackage[english]{babel}
\usepackage{iftex}
\ifPDFTeX
\usepackage[T1]{fontenc}
\usepackage[utf8]{inputenc}
\usepackage{lmodern}
\else
\usepackage{fontspec}
\defaultfontfeatures{Ligatures=TeX}
\setmainfont{Latin Modern Roman}
\setsansfont{Latin Modern Sans}
\setmonofont{Latin Modern Mono}
\usepackage{newunicodechar}
\newfontfamily\unicodeSymbolFont{Segoe UI Symbol}[
  BoldFont=Segoe UI Symbol,
  ItalicFont=Segoe UI Symbol,
  BoldItalicFont=Segoe UI Symbol
]
\newunicodechar{⟂}{{\unicodeSymbolFont ⟂}}
\fi
\usepackage{amsmath,amsfonts}
\usepackage{booktabs,multirow,xcolor,colortbl}
\usepackage{bm}%to use a more powerful version of \boldsymbol
\usepackage{fvextra}
\usepackage{csquotes}
\usepackage{textcomp}
\ifPDFTeX
\DeclareUnicodeCharacter{27C2}{\ensuremath{\perp}}
\fi

\newcommand{\AllSample}{ \mathcal Y^T }

\usepackage[citestyle=authoryear-comp,bibstyle=../Preamble/mystyle_lowercase,maxbibnames=5,minbibnames=1,maxnames=2, minnames=1,datezeros=false,date=long,isbn=false,url=false,backend=biber,useprefix=true]{biblatex}
\addglobalbib{../Preamble/Macro_database_biblatex.bib}	

\usepackage{upquote}
\usepackage{ragged2e}
\usepackage{microtype}

\usepackage{multimedia}
\usepackage{graphicx}
\usepackage[export]{adjustbox}
\usepackage{epstopdf}
\usepackage[labelfont=bf,format=hang, figurename=Fig.]{caption}
\usepackage[subrefformat=parens,labelformat=empty]{subcaption}
\usepackage{pgfpages} % for multiple-screen presentations

\usepackage{array}

\newcolumntype{L}[1]{>{\RaggedRight\let\newline\\\arraybackslash\hspace{0pt}}m{#1}}
\newcolumntype{C}[1]{>{\centering\let\newline\\\arraybackslash\hspace{0pt}}m{#1}}
\newcolumntype{R}[1]{>{\raggedleft\let\newline\\\arraybackslash\hspace{0pt}}m{#1}}
%https://tex.stackexchange.com/questions/12703/how-to-create-fixed-width-table-columns-with-text-raggedright-centered-raggedlef


%%% use serif font for math
\usefonttheme[onlymath]{serif}

%%% define blue highlighting
\newcommand{\highlight}[1]{{\color{blue}{#1}}}

%%%% define command for setting sources

\newcommand{\Source}{}

%%%
\beamertemplatenavigationsymbolsempty


%%% define themes
\usetheme{Boadilla}
\usecolortheme{rose}
\usecolortheme[named=blue]{structure}

\newcommand\textbox[1]{%
  \parbox{.333\textwidth}{#1}%
}

\setbeamertemplate{sections/subsections in toc}[sections numbered]


%%%% set coloring in  foot and head
\setbeamercolor{alerted text}{fg=red!80!black}
\setbeamertemplate{enumerate items}[default] % numbers enumerate items without encircling them
\setbeamertemplate{itemize items}[ball] % numbers enumerate items without encircling them
\setbeamercolor{section in head/foot}{fg=black, bg=white}
\setbeamercolor{subsection in head/foot}{fg=black, bg=white}
\setbeamercolor{title}{fg=black}
\setbeamercolor{subtitle}{fg=blue}

\setbeamertemplate{caption}[numbered]

%%% set overlays as default
\beamerdefaultoverlayspecification{<+->}

%%%% settings for output mode

\mode<beamer>
{
%\setbeameroption{show notes on second screen=right}
}

\mode<trans>
{\renewcommand{\highlight}[1]{{\textbf{#1}}}}


\mode<handout>
{
\renewcommand{\highlight}[1]{{\textbf{#1}}}
%\pgfpagesuselayout{4 on 1}[a4paper,border shrink = 5mm,landscape]
%\setbeamercolor{background canvas}{bg=black!5}
}


%%% define Exercise counter
\newcounter{Exercise}
\resetcounteronoverlays{Exercise}
\setcounter{Exercise}{0}
\newcommand{\Exercise}{\refstepcounter{Exercise}\emph{\structure{Übung \theExercise}: }}

%%%% define environments for theorems etc and make them numbered
\setbeamertemplate{theorems}[numbered]

\newtheorem{assum}{Assumption}
\setbeamertemplate{assum}[numbered]

\newtheorem{defin}{Definition}
\setbeamertemplate{defin}[numbered]

\newtheorem{propos}{Proposition}
\setbeamertemplate{propos}[numbered]


%%%% Make Parts show up in TOC
\makeatletter
\AtBeginPart{%
  \addtocontents{toc}{\protect\beamer@partintoc{\the\c@part}{\beamer@partnameshort}{\the\c@page}}%
}
%% number, shortname, page.
\providecommand\beamer@partintoc[3]{%
  \ifnum\c@tocdepth=-1\relax
    % requesting onlyparts.
    \makebox[6em]{#1.} #2
    \par
  \fi
}
\define@key{beamertoc}{onlyparts}[]{%
  \c@tocdepth=-1\relax
}
\makeatother%

%%%% fix bug in beamer package with note itemize
%%% http://tex.stackexchange.com/questions/1480/changing-the-textwidth-of-the-notes-in-beamer-repost-from-so
\makeatletter
\defbeamertemplate{note page}{infolines}
{%
  \vskip.75em
  \nointerlineskip
  \begin{minipage}{\textwidth} % this is an addition
  \footnotesize \insertnote
  \end{minipage}               % this is an addition
}
\makeatother

\setbeamertemplate{note page}[infolines]


\newcommand{\backupbegin}{
   \newcounter{framenumberappendix}
   \setcounter{framenumberappendix}{\value{framenumber}}
}
\newcommand{\backupend}{
   \addtocounter{framenumberappendix}{-\value{framenumber}}
   \addtocounter{framenumber}{\value{framenumberappendix}}
}

\hypersetup{bookmarksnumbered=true,naturalnames=true,pdfhighlight=/N,citebordercolor={1 1
1},linkbordercolor={1 1 1},colorlinks=true,anchorcolor=black,linkcolor=blue,citecolor=green,urlcolor=green,
pdfpagemode=UseOutlines,breaklinks=true,
pdfkeywords=Dynare, pdfsubject=Dynare,pdfstartpage=1} %pdfpagemode=fullpage

\usepackage[copyright]{ccicons}

\usepackage{tikz,pgfplots}
\pgfplotsset{compat=1.18}
\usetikzlibrary{arrows,calc,decorations.pathreplacing,decorations.shapes}
\pgfdeclarelayer{background}
\pgfdeclarelayer{foreground}
\pgfsetlayers{background,main,foreground}

\usepackage{relsize}
\newcommand\LM{\ensuremath{\mathit{LM}}}
\newcommand\IS{\ensuremath{\mathit{IS}}}
\DeclareMathOperator{\std}{std}


\setbeamertemplate{footline}{
\begin{beamercolorbox}{subsection in head/foot}
\insertsectionnavigationhorizontal{.95\textwidth}{\hskip0pt
\hfill}{\hfill \insertframenumber /\inserttotalframenumber}
\end{beamercolorbox}
}

\newenvironment{witemize}{\itemize\addtolength{\itemsep}{8pt}}{\enditemize}
\newenvironment{wenumerate}{\enumerate\addtolength{\itemsep}{8pt}}{\endenumerate}

\usepackage[]{minted}
\usepackage{tcolorbox}
\usepackage{etoolbox}
\BeforeBeginEnvironment{minted}{\begin{tcolorbox}[boxsep=0mm]}%
\AfterEndEnvironment{minted}{\end{tcolorbox}}%

\AtBeginEnvironment{minted}{%
  \renewcommand{\fcolorbox}[4][]{#4}}  

\usepackage{smartdiagram}

\usepackage{upquote}

\AtBeginSection[]
{
  \begin{frame}{Outline}
    \tableofcontents[currentsection]
  \end{frame}
}


\title[]{\textbf{Advanced Dynare\\}
\medskip{Chapter I}\\
\emph{The software Dynare}
}

\author{Johannes Pfeifer}
\input{../date_place.txt}


\begin{document}

\begin{frame}[plain]
  \titlepage
\end{frame}

\section{Intro}
\begin{frame}{Licensing and Thank you}
	\begin{witemize}
        \item The slides have benefited from the work and materials generously shared by friends and colleagues
        \begin{itemize}
          \item Stéphane Adjemian
          \item Houtan Bastani
          \item Michel Juillard
          \item Willi Mutschler
          \item Philipp Pfeiffer
          \item Marco Ratto
          \item Normann Rion
          \item Sébastien Villemot
        \end{itemize}
        \item The licensing varies between completely free, MIT, CC-BY-SA 4.0, and unspecified\\
        $\rightarrow$ impossible to sort out in time
        \item Where applicable, I state the license\\
        $\rightarrow$ this chapter benefitted from Philipp Pfeiffer
    \end{witemize}
\end{frame}


\begin{frame}{Licensing and Thank you}
	\begin{witemize}
        \item Slides always contain errors
        \item If you spot one, please let me know
        \item I will \alert{correct them in red}
    \end{witemize}
\end{frame}

\section{Dynare Overview}

\begin{frame}{What is Dynare?}
	\begin{witemize}
	    \item A free software platform for handling a wide class of economic models, in particular, DSGE models
        \item Provides modeling language to represent models and computing tasks, with clear separation between
        \begin{enumerate}
          \item particular model declared by the user\\
          $\rightarrow$ binary preprocessor to translate modeling language into model representation useful for computations
          \item computing functions that handle an entire family of DSGE models\\
          $\rightarrow$ collection of MATLAB/Octave routines for computations/task
        \end{enumerate}
	    \item Dynare is very user-friendly and has low entry costs\\
        $\rightarrow$ history is described in the aptly titled \textcite{CherrierEtAl2023}: \citetitle{CherrierEtAl2023}
	    \item Users provide \texttt{.mod} file containing parameters, model equations linking the variables, and computing tasks to be performed
	\end{witemize}
\end{frame}

\begin{frame}{Dynare resources}
	\begin{witemize}
	    \item Download at \url{https://www.dynare.org/}, latest version: Dynare 6.4\\
        $\rightarrow$ Dynare 7 hopefully in Q1 of 2026
        \item Dynare manual: \textcite{Adjemianetal2024}
        \item Forum at \url{https://forum.dynare.org/}
        \item Git repository with code, Wiki, buglist at \url{https://git.dynare.org/Dynare/dynare/}
        \item \textcite{Pfeifer2013}: \href{https://drive.google.com/file/d/1r89OU5OE3CBa6tOlj6l3hNVWEaRH5Anv/view}{\citetitle{Pfeifer2013}}
        \item \textcite{Pfeifer2013Graphs}: \href{https://drive.google.com/file/d/1DO1OK5AKLWfcJX49ZPCGo58Ni2LdERw4/view}{\citetitle{Pfeifer2013Graphs}}
        \item Collection of example files: \url{https://github.com/johannespfeifer/dsge_mod}
	\end{witemize}
\end{frame}


\begin{frame}{Tools available in Dynare}
	\begin{enumerate}
	\item computes the solution of perfect foresight models (arbitrary accuracy)
	\item computes first- and higher-order approximation to solution of stochastic models
	\item estimates (maximum likelihood, Bayesian approach, method of moments) parameters
	of DSGE models, for linear and non-linear models
    \item handles occasionally binding constraints
	\item checks for identification of estimated parameters
	\item computes optimal policy (OSR, Ramsey, discretion)
	\item performs global sensitivity analysis of a model
	\item provides macro language and reporting facility
    \item forecasting
\end{enumerate}
\end{frame}

\begin{frame}{How Dynare works}
\includegraphics[width=0.8\textwidth]{figures/Dynare.PNG}\\
\end{frame}

\section{Modeling language}

\begin{frame}{A Dynare file}
Our model:
\begin{itemize}
    \item Endogenous variables: $\pi_t, y_t, i_t, \varepsilon_t^a, \varepsilon_t^c, \varepsilon_t^i $
    \item Some parameters: $ \beta, \kappa, \sigma,\phi_\pi,\phi_y,\rho^a, \rho^c, \rho^i $
    \item Exogenous variables (shocks): $,u_t^a,u_t^c,u_t^i$
    \item A three equation New Keynesian model:
    \begin{align}
        \pi_t =& \beta \mathbb{E}_t\left[\pi_{t+1}\right]+\kappa y_t +\varepsilon_t^a \\
        y_t =& \mathbb{E}_t \left[y_{t+1} \right]- \frac{1}{\sigma}\left(i_t-\mathbb{E}_t\left[\pi_{t+1}\right]\right)+\varepsilon_t^c \\
        i_t =& \phi_\pi \pi_t + \phi_y y_t+\varepsilon_t^i
    \end{align}
    and AR-processes for exogenous processes
    \begin{equation}
      \varepsilon^j_t=\rho_j \varepsilon_{t-1}^j+u_t^j
    \end{equation}
\end{itemize}
\end{frame}

\begin{frame}{Structure of a Dynare file}
\begin{figure}
\smartdiagramset{back arrow disabled,module minimum width=2.cm,font=\footnotesize}
\smartdiagram[flow diagram:horizontal]{Preamble, Model definition, Steady state, Simulation and estimation}
\end{figure}
\begin{enumerate}
\item Preamble: Declaration of endogenous and exogenous variables and parameters. Set value for calibrated parameters
\item Model definition: State the dynamic model equations
\item Steady state (long-run equilibrium)
\item Additional computations: model checks,  simulation, estimation, ...
\end{enumerate}
\begin{itemize}
\item All command lines are terminated with a semicolon
\item Blocks are terminated with an \texttt{end;}
\item Comments either using \% (Matlab style) or //, /* */ (C-style)\\
$\rightarrow$ Matlab comments allow simpler debugging
\end{itemize}
\end{frame}

\begin{frame}[fragile]{How to interact with Dynare: 1) Preamble}
\begin{columns}<.->
\begin{column}<1->[T]{0.5\textwidth}
	    \begin{witemize}
	        \item Endogenous variables: $\pi_t, y_t, i_t, \varepsilon_t^a, \varepsilon_t^c, \varepsilon_t^i$
	        \item Exogenous variables (shocks): $u_t^a, u_t^c, u_t^i$
	        \item Parameters: $ \beta, \theta, \eta, \sigma,\phi_\pi,\phi_y,\rho^a,\rho^c,\rho^i $
	    \end{witemize}
\end{column}
\begin{column}[T]{0.5\textwidth}
Dynare code
\begin{minted}{c++}
var // endo vars
pie y inom eps_a eps_c eps_i;

varexo
u_a u_c u_i;

parameters
betta theta eta sig phi_pi
phi_y rho_a rho_c rho_i;
\end{minted}
\end{column}
\end{columns}
\end{frame}

\begin{frame}[fragile]{\LaTeX\ definitions in Dynare}
\begin{itemize}
	\item Note: For all three types, you can specify \LaTeX\ names with \$\$ and long names for Dynare output with \texttt{(long\_name='name here')}
\end{itemize}

\begin{minted}{c++}
var // endo vars
pie ${\pi}$ (long_name='Quarterly inflation rate')
y   ${y}$   (long_name='Output gap')
;

varexo
u_a ${u_a}$ (long_name='Cost-push shock')
;

parameters
betta   ${\beta}$    (long_name='Discount factor')
theta   ${\theta}$   (long_name='Calvo parameter')
;
\end{minted}
\end{frame}


\begin{frame}[fragile]{How to interact with Dynare: 1) Preamble (cont'd)}
\begin{columns}<.->[T] % align columns
\begin{column}<1->[T]{0.5\textwidth}
	    \begin{witemize}
	        \item Set $\beta$ and $\theta$ to conventional values informed by real interest rates and price adjustment frequencies
	        \item Choose parameters for the Taylor rule
	        \item Set autocorrelation coefficients
	        \item Unit intertemporal elasticity of substitution
	    \end{witemize}
\end{column}
\begin{column}[T]{0.5\textwidth}
\begin{minted}{c++}
// Calibration
betta   = 0.995;
eta     = 1;
theta   = 0.78;
phi_pi  = 1.5;
phi_y   = .5/4;
rho_a   = 0.8;
rho_c   = 0.9;
rho_i   = 0.5;
sig     = 1;
\end{minted}
\end{column}
\end{columns}
\end{frame}

\begin{frame}[fragile]{How to interact with Dynare: 2) Model definition}
\begin{columns}<.->[T] % align columns
\begin{column}<1->[T]{0.4\textwidth}
	    \begin{witemize}
	        \item Dynare allows writing the model in a very natural way
	        \item \texttt{model}-block contains model declaration\\
            $\rightarrow$ \texttt{linear} states that model is already linearised
	        \item \texttt{(+1)} indicates one-period ahead variables $X_{t+1}$ (forward-looking)
	        \item \texttt{(-1)} indicates lagged variables $X_{t-1}$ (backward-looking)
            \item $X_t$ is written $X$ (no time index needed for the current period)
	    \end{witemize}
\end{column}
\begin{column}[T]{0.6\textwidth}
\begin{minted}{c++}
model(linear);
[name='Phillips Curve']
pie=betta*pie(+1)+kappa*y+zeps_a;
[name='DIS curve']
y=y(+1)-1/sig*(inom-pie(+1))
+zeps_c;
[name='Nominal interest rate Taylor rule']
inom=phi_pi*pie+phi_y*y+zeps_i;
zeps_a=rho_a*zeps_a(-1)+u_a;
zeps_c=rho_c*zeps_c(-1)+u_c;
zeps_i=rho_i*zeps_i(-1)+u_i;

end;
\end{minted}
\end{column}
\end{columns}
\end{frame}

\begin{frame}{Timing conventions}
\begin{itemize}
	\item Time indices refer to period when variable is determined, i.e.\ \highlight{stock at the end of period notation}
	\item Capital stock used in production at time $t$ is determined by investment at time $t-1$\\
$\rightarrow$ predetermined
	\begin{align}
	K_t&=(1-\delta)K_{t-1}+I_t\\
	Y_t&=K_{t-1}^\alpha L_t^{1-\alpha}
	\end{align}
	\item To use \highlight{stock at the beginning of period notation}
	\begin{align}
	K_{t+1}&=(1-\delta)K_{t}+I_t\\
	Y_t&=K_{t}^\alpha L_t^{1-\alpha}
	\end{align}
	you can declare variable as predetermined using \texttt{predetermined\_variables} declaration
\end{itemize}
\end{frame}


\begin{frame}[fragile]{Timing conventions: Code}
\begin{columns}<.->
\begin{column}<1->[T]{0.5\textwidth}
Example using the default convention
\begin{minted}{c++}
var y, k, i;
...
...
model;
y = k(-1)^alpha;
k = i + (1-delta)*k(-1);
...
end;
\end{minted}
\end{column}
\begin{column}[T]{0.5\textwidth}
Using \href{https://www.dynare.org/manual/the-model-file.html?highlight=predetermined#predetermined_variables}{\texttt{predetermined\_variables}}:
\begin{minted}{c++}
var y, k, i;
predetermined_variables k;
...
model;
y = k^alpha;
k(+1) = i + (1-delta)*k;
...
end;
\end{minted}
\end{column}
\end{columns}
\end{frame}

\begin{frame}{Timing conventions: conditional expectations}
\begin{witemize}
\item Conditional expectations operator $\mathbb{E}_t$ typically not needed in Dynare as it starts from equilibrium system
\begin{equation}
0=\mathbb{E}_t\left(f(y_{t+1},y_t,y_{t-1},\varepsilon_t)\right)
\end{equation}
and therefore assumes conditional expectations operator in front of every equation
\item \texttt{1/c=beta/c(+1)*(R(+1));}\\
corresponds to
\begin{equation}
\frac{1}{c_t}=\mathbb{E}_t\left(\frac{1}{c_t}\right)=\mathbb{E}_t\left(\beta \frac{R_{t+1}}{c_{t+1}}\right)=\beta \mathbb{E}_t\left(\frac{R_{t+1}}{c_{t+1}}\right)
\end{equation}
\item Due to certainty equivalence at first order, we do not have to bother with this, but intricate issues can arise at higher order (e.g.\ Epstein-Weil-Zin preferences)
    \item \texttt{EXPECTATION(INTEGER)(MODEL\_EXPRESSION)} allows specifying information set of expectations explicitly
\end{witemize}
\end{frame}


\begin{frame}{The steady state}
\begin{witemize}
	\item If not told differently, Dynare tries whether 0 is the steady state, and if not, starts a numerical solver to find the steady state
	\item As 0 is often infeasible for nonlinear models, always provide explicit steady state information
	\item Three ways to do this:
	\begin{enumerate}
		\item compute steady state analytically and provide it to Dynare using \texttt{steady\_state\_model}-block
		\item provide initial guesses to start the numerical optimizer using the \texttt{initval}-block (the better, the higher the chances of finding the steady state; economic intuition provides good guidance, e.g. $I<C<Y<K$, with $I=\delta K$.)
		\item use explicit steady state file, which offers a lot of flexibility
	\end{enumerate}
	\item Note: linearized models typically do not spare you from computing the steady state, they just shift the problem to computing parameters
\end{witemize}
\end{frame}


\begin{frame}[fragile]{How to interact with Dynare: 3) Steady state}
\begin{columns}<.->
\begin{column}<1->[T]{0.5\textwidth}
	    \begin{witemize}
	        \item Our model is expressed in deviation from the steady state (around zero inflation)
	        \item The deterministic steady state implies no shocks: $u= 0$
	        \item Thus, all variables are zero in the steady state
	    \end{witemize}
\end{column}
\begin{column}[T]{0.5\textwidth}
\begin{minted}{c++}
steady_state_model;
pie=0;
y=0;
inom=0;
zeps_a=0;
zeps_c=0;
zeps_i=0;
end;
\end{minted}
or
\begin{minted}{c++}
steady_state_model;
end;
\end{minted}
\end{column}
\end{columns}
\end{frame}


\begin{frame}[fragile]{Numerical steady state}
\begin{witemize}
    \item Dynare can solve for the steady state using initial starting values provided
\begin{minted}{c++}
initval;
y=0;
end;
\end{minted}
    \item Inside \texttt{steady\_state\_model;} and \texttt{initval;}-blocks, you can define variables recursively and use auxiliary definitions
    \item Always use a fully analytical steady state if possible (especially for \texttt{estimation})
    \item The steady state values stored in \texttt{oo\_.steady\_state} follow the declaration order
    \item You can access steady state values within the \texttt{model;}-block using the \texttt{STEADY\_STATE (MODEL\_EXPRESSION)} operator
\end{witemize}
\end{frame}

\begin{frame}[fragile]{Steady state \texttt{.m}-file}
\begin{witemize}
    \item You can also calculate the steady state using your own numerical routines.
    \item Create a Matlab file with the same name as your model followed by \texttt{\_steadystate.m}\\
    $\rightarrow$ \texttt{DSGE\_steadystate.m} for model file \texttt{DSGE.mod}
    \item The file allows a lot of flexibility like calling numerical solvers, but syntax is more involved
    \item Full example available \href{https://git.dynare.org/Dynare/dynare/-/blob/master/examples/NK_baseline_steadystate.m}{here [LINK]}
\end{witemize}
\end{frame}

\begin{frame}{Digression: Log-linearization vs.\ linearization}
\begin{witemize}
\item Dynare linearizes the model around the steady state\vspace{0.2cm}
\item However, we often want log-linearization to express everything in percentage deviations from steady state (or the balanced growth path)\vspace{0.2cm}
\item To obtain a log-linear solution, you can\vspace{0.2cm}
\begin{enumerate}
\item use the \texttt{loglinear} option\vspace{0.2cm}
\item write down the model in terms of logarithmic transformations of the original variables (exp()-substitution)\vspace{0.2cm}
\item append the logarithm of model variables as separate model variables (e.g.\ \texttt{log\_y=log(y)})\\

$\rightarrow$ Dynare 6.x: use \texttt{var(log) y;} to append \texttt{LOG\_Y}
\end{enumerate}
\item The last option tends to be the most preferable
\end{witemize}
\end{frame}

\begin{frame}[fragile]{Structure of typical Dynare $.mod$ file: shocks}
\begin{witemize}
	\item \texttt{shocks;}-block defining the covariance matrix of exogenous shocks
\begin{minted}{c++}
shocks;
    var u_c  = 0.005^2;
    var u_a  = 0.005^2;
    var u_i  = 0.005^2;
    //corr u_c,u_i=0.5;
end;
\end{minted}
\item Also allows specifying correlations (or covariances)
\item Different syntax for perfect foresight and OccBin
\end{witemize}
\end{frame}

\begin{frame}[fragile]{Parameter dependence}
\begin{witemize}
    \item Sometimes, parameters like the slope of the Phillips Curve are functions of other parameters:
    \begin{equation}
      \kappa=\frac{(\sigma+\eta)(1-\theta)(1-\beta\theta)}{\theta}
    \end{equation}
    \item There are two ways to handle such objects:
    \begin{enumerate}
        \item As a model-local variable in the \texttt{model} block, using the \# operator\\
        $\rightarrow$ declaring \texttt{model\_local\_variable} allows setting \LaTeX\ name
        \item setting them in a \texttt{steady\_state\_model}-block or file
    \end{enumerate}

\begin{minted}{c++}
model_local_variable kappa $\kappa$;
model(linear);
#kappa = (sig+eta)*(1-theta)*(1-betta*theta)/theta;

pie=betta*pie(+1)+kappa*y+zeps_a;
\end{minted}
    \item Accounting for these dependencies is important, especially for estimation
\end{witemize}
\end{frame}


\begin{frame}[fragile]{Requesting computations}
\begin{witemize}
	  \item Display residuals with \href{https://www.dynare.org/manual/the-model-file.html?highlight=resid#resid}{\texttt{resid;}} command
	  \item Compute and display steady state with \href{https://www.dynare.org/manual/the-model-file.html?highlight=steady#steady}{\texttt{steady;}} command
	\item \href{https://www.dynare.org/manual/the-model-file.html?highlight=check#check}{\texttt{check;}} command to display local stability conditions
\begin{minted}{c++}
resid;  %display residuals
steady; %display steady state (and compute it if necessary)
check;  %display local stability information
\end{minted}
	\item \href{https://www.dynare.org/manual/the-model-file.html?highlight=stoch_simul#stoch_simul}{\texttt{stoch\_simul;}} command  to
	\begin{itemize}
		\item linearize the model around the steady state
		\item compute recursive equilibrium laws of motion of the linearized model
		\item analyze the model via impulse responses, stochastic simulations, moments, etc.
	\end{itemize}\vspace{0.2cm}
\begin{minted}{c++}
stoch_simul(order=1,irf=20,irf_plot_threshold=0) y pie inom;
\end{minted}
\end{witemize}
\end{frame}

\begin{frame}[fragile]{\LaTeX\ support}
\begin{witemize}
    \item 	 Commands to generate \LaTeX\ output
\begin{minted}{c++}
write_latex_dynamic_model(write_equation_tags);
write_latex_parameter_table;
write_latex_definitions;
collect_latex_files;
\end{minted}
\item Particularly useful for debugging
\end{witemize}
\end{frame}

\section{Invocation and results}

\begin{frame}[fragile]{Running Dynare}
\begin{witemize}
    \item You need to add Dynare to your Matlab path, e.g.
    \begin{minted}{matlab}
    addpath('c:/dynare/6.4/matlab')
    \end{minted}
    \item You can run a \texttt{.mod}-file located in MATLAB's \textit{Current Folder} using
    \begin{minted}{matlab}
    dynare FILENAME[.mod]
    \end{minted}
\end{witemize}
\end{frame}


\begin{frame}[fragile]{Accessing Dynare output}
\begin{witemize}
    \item Dynare stores results, model information, and selected options in MATLAB structures (use the tab key at the MATLAB command line to access the different fields)
    \begin{witemize}
        \item \texttt{oo\_} stores results from computations (simulation results, ...)
        \item \texttt{M\_} stores model information (e.g. info on parameters and their values)
        \item \texttt{options\_} stores the selected options
        \item \texttt{estim\_params\_} stores information on estimated parameters
        \item \texttt{bayestopt\_} stores information on the prior
    \end{witemize}
    \item All output fields should be documented in the manual
\end{witemize}
\end{frame}



\part{References}
%%%
\begin{frame}[allowframebreaks]{Bibliography}
\printbibliography
\end{frame}
\end{document}
